\documentclass[11pt,a4paper]{article}

\usepackage[margin=1in]{geometry}
\usepackage{times}
\usepackage{setspace}
\usepackage{microtype}
\usepackage[hidelinks]{hyperref}

\onehalfspacing
\setlength{\parskip}{0.4em}
\setlength{\parindent}{1.2em}

\title{Biphasic Pathogenesis of \textit{Yersinia pestis}:\\
Intracellular Conditioning and Extracellular Resistance to Re-phagocytosis}
\author{}
\date{}

\begin{document}
\maketitle
\thispagestyle{empty}

\section{Introduction}

\textit{Yersinia pestis}, the causative agent of plague, is a facultative intracellular Gram-negative bacterium whose disease course depends on an orderly transition between two physiological states. In the first phase the organism is taken up by host phagocytes---especially macrophages---where it survives and multiplies while reprogramming its surface and virulence repertoire. In the second phase it exits those cells as an extracellular pathogen that resists re-uptake by neutrophils and macrophages and multiplies freely in tissue and blood. This biphasic strategy is central to both bubonic plague, classically initiated by flea bite, and pneumonic plague, initiated by inhalation of infectious aerosols. The sections that follow first outline the two infection routes, then what differs between the bacterial phases and which environmental and regulatory cues drive the switch, and finally why natural transmission often forces the sequence rather than allowing the organism to begin already armed.

\section{Pathogenesis of bubonic and pneumonic plague}

In bubonic plague, \textit{Y.\ pestis} is transmitted from an infected flea. In the flea midgut the bacterium grows at roughly ambient temperature (about $26$--$28^\circ\mathrm{C}$) and lacks the full mammalian-temperature virulence package, including the proteinaceous F1 capsule. After dermal inoculation, most bacilli that encounter neutrophils are killed. A minority are phagocytosed by tissue macrophages. Macrophages fail to sterilise the infection; instead they provide a protected niche in which the bacteria replicate and ``arm'' themselves by inducing virulence factors that will be needed outside the cell. From this local reservoir the organism spreads via lymphatics to regional lymph nodes, producing the characteristic swollen and necrotic buboes. Continued extracellular proliferation can seed the bloodstream (secondary septicaemic plague) and, in some cases, the lungs (secondary pneumonic plague).

Primary pneumonic plague follows inhalation of infectious droplets. Early in pulmonary infection, alveolar macrophages are among the first host cells contacted; neutrophils are also engaged and become more prominent as disease progresses. As in the dermal route, an initial intracellular or closely cell-associated phase is followed by massive extracellular outgrowth in the airways and alveolar spaces. Clinically and experimentally, primary pneumonic plague is often described as biphasic at the host level as well: an early phase with relatively limited overt inflammation is succeeded by a rapidly progressive, highly inflammatory necrotising pneumonia in which extracellular bacilli dominate tissue sections. Autopsy and experimental material from advanced pneumonic plague typically show abundant free bacteria with little evidence of ongoing phagocytosis---consistent with successful transition into the anti-phagocytic state.

Both routes therefore share a sequential pathogenesis: limited early association with (and survival inside) phagocytes, followed by release and extracellular dominance once anti-phagocytic systems are in place.

\section{Phase one: survival and replication in macrophages}

The first phase is defined by entry into, and adaptation within, professional phagocytes. \textit{Y.\ pestis} lacks functional Invasin and YadA---adhesins important in the related enteropathogenic yersiniae---yet still associates with and enters macrophages and certain non-professional phagocytes. Surface proteins implicated in early attachment and uptake include the outer-membrane protease and adhesin Pla (plasminogen activator, encoded on plasmid pPCP1), the adhesin Ail, and additional factors such as YadBC. These molecules promote contact with host cells and, in some experimental systems, promote internalisation.

Once inside macrophages, the bacteria occupy a specialised compartment often termed the \textit{Yersinia}-containing vacuole (YCV). Rather than being delivered efficiently to a fully bactericidal phagolysosome, the YCV can acquire late endosomal or lysosomal markers while remaining spacious and relatively permissive; acidification and autophagic killing are limited. Within this niche the organism multiplies and experiences a suite of host-associated stresses: mammalian body temperature ($37^\circ\mathrm{C}$), reduced pH, osmotic challenge, reactive oxygen species, and altered metal availability. These conditions do more than select for hardy cells; they act as signals that rewire gene expression. Intracellular survival also buys time for the synthesis of the anti-phagocytic machinery that will define the second phase. Neutrophils, by contrast, are generally more bactericidal for \textit{Y.\ pestis} early on, so the macrophage pathway is thought to be the more productive route for establishing a conditioned, dissemination-competent population.

Bacteria arriving from the flea (or grown at flea-like temperature in vitro) remain relatively susceptible to phagocytosis and killing. The intracellular interval is therefore a developmental window as much as a refuge: under host temperature and phagosomal stress the population becomes equipped for life outside the cell.

\section{Phase two: surface remodelling and resistance to re-phagocytosis}

After escape from macrophages---by mechanisms that remain incompletely defined but that may involve host-cell death pathways observed in vitro---the matured bacteria face neutrophils and macrophages in the extracellular space. Survival now depends on active resistance to re-phagocytosis and on suppression of inflammatory responses that would otherwise clear the infection.

Three classes of factor are especially important. First, the F1 (Caf1) capsule, assembled as linear fibres of Caf1 subunits, is produced in large amounts at $35$--$37^\circ\mathrm{C}$ and coats the bacterial surface. F1 physically interferes with uptake and can mask shorter outer-membrane adhesins. Second, the pH~6 antigen (Psa), a fimbrial adhesin induced at $37^\circ\mathrm{C}$ and under mildly acidic conditions typical of phagosomal environments, promotes attachment in a way that does not efficiently trigger internalisation and can further inhibit phagocytosis, in part by binding host lipoproteins that shield the bacterium. Third, and central to virulence, is the plasmid-encoded type III secretion system (T3SS; Ysc--Yop). Upon contact with host cells, this needle-like apparatus injects effector proteins (Yops) into the eukaryotic cytosol. YopH, a potent protein tyrosine phosphatase, disrupts focal-adhesion signalling required for phagocytic uptake. YopE and YopT inactivate Rho-family GTPases that organise the actin cytoskeleton; YpkA/YopO further disables Rho signalling and G-protein pathways that support cytoskeletal remodelling. Together these effectors paralyse phagocytosis and blunt pro-inflammatory signalling. LcrV (V antigen), associated with the T3SS tip, also contributes to immune modulation.

In the matured state, adhesive structures such as Ail and Pla may still be present, but the combination of a dense F1 coat, Psa, and T3SS effectors shifts the balance from ``enterable'' to ``unengulfable.'' The result is predominantly extracellular growth in lymph nodes, blood, and lung---the phase associated with fulminant tissue destruction and high bacterial burden.

\section{What differs between the phases, and what causes the switch}

The two phases differ in lifestyle (intracellular or closely cell-associated versus extracellular), in surface architecture (invasion-competent adhesins relatively exposed versus capsule- and Psa-dominated, T3SS-armed surfaces), and in interaction with innate immunity (exploitation of macrophage permissiveness versus active blockade of neutrophil and macrophage uptake).

The switch is driven largely by environmental cues that report transition from the vector (or ambient) niche into the mammalian host and then into the phagosomal microenvironment. Temperature is the master signal. A shift from about $26$--$28^\circ\mathrm{C}$ to $37^\circ\mathrm{C}$ induces expression of F1, components of the T3SS, and many other virulence loci. Thermoregulation of the Ysc--Yop system involves the transcriptional activator LcrF (VirF); at lower temperature, RNA thermosensing and the histone-like repressor YmoA restrain expression, whereas host temperature relieves that repression and allows build-up of the secretion apparatus. Full deployment of Yop injection further requires host-cell contact and is classically linked to the low-calcium response in vitro. Mildly acidic pH, as encountered in the phagosome or YCV, specifically favours induction of Psa (pH~6 antigen), coupling intracellular residence to expression of an anti-phagocytic adhesin. Additional stresses inside the macrophage---oxidative pressure, metal limitation, and other phagosomal signals---reinforce transcriptional reprogramming through regulators such as PhoP and RovA that affect intracellular survival and selected virulence traits.

A structural consequence of this reprogramming completes the phenotypic switch. As F1 polymerises over the cell surface, short adhesins such as Ail and Pla become sterically less able to engage host receptors that would promote uptake, whereas longer structures such as Psa can still project outward but channel the interaction toward adhesion without productive internalisation and toward efficient T3SS delivery. Thus the same cell that earlier needed invasive factors to enter macrophages later presents a surface optimised to refuse re-entry.

Several mechanistic details remain open: precisely how bacteria exit the YCV and the macrophage in vivo; how the timing of F1, Psa, and T3SS induction is coordinated so that escape and anti-phagocytic competence coincide; and how route-specific cues in skin versus lung fine-tune the same biphasic programme. Across experimental systems the surrounding model is nonetheless consistent: early macrophage infection conditions \textit{Y.\ pestis} at host temperature and under phagosomal stress; induction of capsule, pH~6 antigen, and the Yop T3SS then converts the population into an extracellular pathogen that resists re-phagocytosis and drives the destructive phases of bubonic and pneumonic plague.

\section{Why two phases? Why not begin already in phase two?}

If the second, anti-phagocytic state is what enables explosive extracellular growth, why does \textit{Y.\ pestis} not arrive already in that state and dispense with the intracellular interval? The literature treats the two-phase programme as the joint product of transmission biology, the time needed to reconfigure gene expression, the mutual incompatibility of invasive and anti-phagocytic surface programmes, and the fitness cost of constitutively deploying expensive virulence systems. Under natural conditions the bacterium often \emph{cannot} begin in phase two; even when experimental designs approximate that phenotype, skipping a protected adaptation window remains risky.

\subsection{Transmission from the flea delivers an unarmed population}

The strongest constraint is ecological. In the classical bubonic cycle, \textit{Y.\ pestis} multiplies in the flea at ambient temperature (roughly $20$--$28^\circ\mathrm{C}$). Under those conditions the F1 capsule is produced poorly or not at all, and the suite of temperature-induced mammalian virulence factors---including much of the Ysc--Yop T3SS---is not fully on. After a flea bite the inoculum that enters the dermis is therefore a vector-adapted population, not one already coated in F1 and primed for Yop injection. Vadyvaloo and colleagues showed that transit through the flea does impose a distinctive pre-transmission transcriptional state and can confer some resistance phenotypes, but they also emphasised that many of the classic anti-innate-immunity virulence factors are induced only after the temperature shift to $37^\circ\mathrm{C}$ in the mammalian host. Until that induction has occurred, the bacteria remain comparatively vulnerable to neutrophils and other early defences. Phase one is, in large part, what happens while an unarmed flea-derived inoculum becomes armed.

Laboratory experiments that pre-grow bacteria at $37^\circ\mathrm{C}$ before infection can obscure this constraint. Such pre-incubation can induce F1 and the T3SS and increase resistance to phagocytosis, giving the culture a head start toward phase two. That shortcut is useful for dissecting mechanisms, but it does not recapitulate natural flea transmission. In nature, phase two is unavailable at the moment of inoculation; it must be built after entry into the warm host.

\subsection{Induction takes time; neutrophils do not wait}

Even after the temperature cue is received, reprogramming is not instantaneous. Capsule assembly, build-up of the T3SS apparatus, and expression of Psa and other surface factors require hours of transcription, translation, and assembly. During that lag a small dermal inoculum is exposed to phagocytes. Neutrophils are generally efficient killers of \textit{Y.\ pestis} that have not yet acquired anti-phagocytic competence. Macrophages, by contrast, are more permissive: they take bacteria up into a YCV that supports survival and replication while virulence gene expression catches up. Reviews of early plague pathogenesis therefore cast the macrophage as a sheltered incubator in which the population can expand and complete temperature- and stress-dependent maturation before facing the extracellular innate immune system in force. Without that shelter, stochastic killing of a small, still-adapting inoculum could extinguish the infection before extracellular dominance is achieved.

\subsection{Phase-two surfaces work against phase-one entry}

A second, more molecular reason is that the two phases are partly antagonistic rather than merely sequential versions of one surface. The invasive adhesins that promote early contact and uptake (Pla, Ail, and related factors) act at a short range from the outer membrane. The F1 capsule that characterises the matured anti-phagocytic state forms a fibre coat that can sterically mask those shorter adhesins and reduce productive engagement with host receptors that drive internalisation. Psa and the Yops actively discourage uptake. Ke and co-workers framed this as a regulatory problem: the bacterium must use invasive factors early and anti-phagocytic factors later, and it must switch between programmes that pull in opposite directions. If \textit{Y.\ pestis} constitutively presented a full phase-two surface, it would impair the very macrophage entry that, under natural transmission, provides the protected niche for adaptation. Temporal separation---invade and condition first; refuse re-uptake second---resolves that conflict.

\subsection{Constitutive phase two is metabolically and ecologically costly}

A third line of argument concerns fitness cost and dual-host lifestyle. Full induction of the Ysc--Yop system at mammalian temperature is linked to the classical low-calcium response: under inducing conditions in vitro, unrestricted Yop expression is associated with growth restriction. Maintaining a permanently assembled, secretion-competent T3SS is energetically expensive and is tightly gated by temperature (LcrF/YmoA thermosensing), calcium-related checkpoints, and host-cell contact. Selection therefore favours keeping the system off in environments where it is not needed---notably the flea midgut and the earliest moments after transmission---and turning it on only when host cues justify the cost. The same logic applies more broadly to temperature-regulated mammalian virulence factors: constitutively expressing a mammalian programme would be wasteful or counterproductive in the vector, just as constitutively expressing a flea-adapted programme would leave the bacterium unarmed in the mammal. For a pathogen that must succeed in two very different hosts, biphasic, signal-dependent control is a coherent evolutionary answer.

\subsection{Limits of the model}

Single elements of this account should not be overstated. F1-negative strains can retain substantial virulence in some animal models, indicating that capsule alone does not explain the need for an early intracellular interval; the broader conditioning programme, especially T3SS competence, is more central. Pre-growth at $37^\circ\mathrm{C}$ can partially substitute for in-host induction, which shows that what ultimately matters is the \emph{phenotype} (armed versus unarmed), not the macrophage as an absolute obligate residence. In primary pneumonic plague the donor is already a warm-blooded host, so inhaled bacteria may arrive closer to a phase-two state than flea-regurgitated bacteria do; even so, early targeting of alveolar macrophages and neutrophils is well documented, and disease in the lung remains biphasic at the host-response level. Whether intracellular replication is required in every successful infection is still being refined: some bacteria may complete much of their adaptation extracellularly if inoculum size, pre-adaptation, and local phagocyte dynamics permit. Under natural bubonic transmission, however, phase two is unavailable at the outset, and the first phase remains the adaptation and amplification window forced by vector temperature, induction lag, and neutrophil pressure. The two-phase design is then the strategy that makes plague possible when the bacterium must move from flea to mammal with an initially unarmed surface.

\section{Working hypothesis: temporary macrophage exploitation as a staging strategy}

The account so far establishes that \textit{Y.\ pestis} uses a temporary intracellular (or closely macrophage-associated) interval before dominating as an extracellular pathogen, and that natural transmission often forces that sequence. It does not yet state a positive hypothesis for what functional pay-off the macrophage stage is selected to deliver beyond buying time to arm. We propose a working hypothesis with two route-weighted strands: use of host phagocytes as vehicles and portals into nutrient-rich lymphoid tissue in bubonic disease, and use of an early, relatively protected growth window to accumulate biomass before a delayed wave of harsh neutrophilic inflammation in pneumonic disease. On this view the macrophage is a short-lived staging compartment---a host-derived niche that solves an early-infection bottleneck of small inoculum, incomplete virulence induction, and lethal phagocytes---so that a later extracellular programme can succeed.

\subsection{Statement of the hypothesis}

We take the central claim to be evolutionary rather than intentional: \textit{Yersinia pestis} is selected to accept a transient macrophage-associated phase because that phase multiplies the probability that a limited, initially under-armed inoculum will reach a critical population size and a favourable anatomical site before fully competent extracellular innate immunity---especially neutrophils---clears the infection. Once that threshold is crossed and anti-phagocytic systems (F1, Psa, T3SS) are in place, the organism exits the staging niche and grows as an extracellular pathogen. The form of the bottleneck differs by route. In bubonic plague it is largely spatial and nutritional (getting from dermis to a lymphoid replication site). In pneumonic plague it is largely temporal and immunological (expanding during a pre-inflammatory window before a destructive neutrophil-dominated response).

\subsection{Bubonic plague: carriage and delivery into a lymphoid growth niche}

In flea-borne infection the inoculum is deposited in the dermis, far from the lymph nodes in which the characteristic bubo later forms. A long-standing model holds that tissue macrophages (and related mononuclear phagocytes) take up bacteria at the bite site and, by ordinary trafficking toward draining lymph nodes, deliver them there---a Trojan-horse style of dissemination. Once in the node, bacteria exit or are released from dying host cells and switch to extracellular replication, producing the enormous bacterial burdens associated with buboes and subsequent septicaemia.

Two related pay-offs follow. First, macrophage (or dendritic-cell or monocyte) migration provides a protected conduit from skin to node, reducing free exposure to extracellular effectors and neutrophils during transit. Sphingosine-1-phosphate (S1P)--dependent trafficking of intracellular \textit{Y.\ pestis} through lymphoid tissue has been reported, consistent with active cellular carriage contributing to bubo formation and onward spread. Second, draining lymph nodes concentrate fluid, cells, and metabolites drained from the infected territory. Relative to the sparsely vascularised dermis at a flea bite, the node offers a denser supply of host-derived nutrients (including iron and other metals that \textit{Y.\ pestis} must scavenge) and a structured microenvironment in which extracellular outgrowth, once anti-phagocytic systems are online, can become explosive. The macrophage stage then relocates the infection to a site of denser host-derived nutrients and systemic connectivity, rather than merely hiding bacteria in the skin.

Recent imaging and microneedle studies complicate any exclusive Trojan-horse story: free bacteria can also reach draining nodes within minutes via lymph flow, with little requirement for intracellular carriage. The hypothesis does not require Trojan-horse transport to be the only path to the node. It requires only that macrophage association remain advantageous when it occurs---sheltering a fraction of the inoculum, perhaps supporting slower S1P-linked cell-to-cell or node-to-node movement, and coupling physical delivery to the same intracellular conditioning that produces phase-two competence. Free lymph drainage and cellular carriage can coexist. The prediction is that impairing early mononuclear phagocyte permissiveness or trafficking will reduce the fraction of infections that establish high nodal burdens from small dermal inocula, even if some free bacteria still arrive.

\subsection{Pneumonic plague: early build-up before a neutrophil wave}

Primary pneumonic plague is biphasic at the host level. An early pre-inflammatory phase allows substantial bacterial growth in the lung with limited overt inflammation; this is followed by a late, highly inflammatory phase dominated by neutrophils, tissue necrosis, and massive extracellular bacterial loads. Experimental work shows delayed neutrophil recruitment relative to other pulmonary pathogens, and early preferential interaction of the T3SS with alveolar macrophages (and then neutrophils).

Under this strand of the hypothesis, temporary exploitation of the alveolar macrophage environment---and more broadly of a quiet early lung milieu---serves a different primary purpose from nodal transport. A modest inhaled inoculum that immediately faced a full neutrophilic assault would often be extinguished. By residing in or closely engaging alveolar macrophages, which are more permissive than neutrophils early on, and by delaying inflammatory recruitment (in part through T3SS effectors and related immune modulation), \textit{Y.\ pestis} gains time to increase numbers and complete anti-phagocytic maturation. When neutrophils finally flood the lung, the bacterial population is no longer a handful of adapting cells but a large, largely extracellular, phase-two-competent population able to resist re-phagocytosis and continue expanding despite the inflammatory storm. The early macrophage-associated interval is then a temporal staging step: build biomass and arm before the more lethal phagocyte class arrives in force.

That early protected interval supplies the numerical and phenotypic head start needed for infection to persist. It also explains why pneumonic disease can remain biphasic even when inhaled bacteria come from another warm host and may already express more of the phase-two programme than flea-regurgitated organisms: the bottleneck includes temperature induction, but also the race between bacterial expansion and the timing of neutrophilic inflammation.

\subsection{Route weighting and predictions}

Both strands are instances of the same staging logic. In bubonic disease the early bottleneck is a small dermal inoculum at a distance from the lymphoid growth site; the macrophage stage supplies shelter and trafficking toward the node, with conditioning en route, and exit into extracellular phase-two growth enables explosive replication in a nutrient-rich, systemically connected node. In pneumonic disease the early bottleneck is a small inhaled inoculum facing an impending neutrophil influx; the macrophage stage supplies shelter and early replication with conditioning in the lung, and exit into phase two enables survival and outgrowth through late neutrophilic pneumonia.

Anything that shortens or abolishes the protective staging interval, without providing an equivalent head start, should raise the minimum infectious dose and increase early clearance. Conversely, artificial head starts---large inocula, pre-induction at $37^\circ\mathrm{C}$, experimental delay of neutrophil recruitment---should reduce dependence on prolonged macrophage residence. Route-specific predictions follow. In bubonic models, disrupting mononuclear trafficking or nodal colonisation should matter more than in pneumonic models. In pneumonic models, accelerating early neutrophilic inflammation should collapse the pre-inflammatory growth window and blunt disease more than blocking lymph-node delivery.

\subsection{Scope of the claim}

The claim is not that \textit{Y.\ pestis} is an obligate intracellular pathogen, nor that every successful bacterium must pass through a macrophage. Lymph nodes are not uniquely rich in host nutrients in any absolute sense; relative to the flea-bite dermis they are a superior site for rapid extracellular expansion and systemic seeding once phase two is achieved. Free lymphatic spread is not denied. Temporary macrophage utilisation is selected, on this account, because it systematically improves the odds of clearing early bottlenecks, with transport and placement emphasised in bubonic disease and pre-neutrophil amplification emphasised in pneumonic disease.

\subsection{Empirical tests}

Several tests follow directly. Infection outcomes can be compared for flea-like (low-temperature) versus pre-induced ($37^\circ\mathrm{C}$) inocula at matched small doses, with and without early neutrophil depletion or enhancement. The fraction of dermal inoculum that reaches nodes free versus cell-associated can be tracked, asking whether cell-associated bacteria are enriched for phase-two markers at arrival. In pneumonic models, bacterial population size can be quantified at the transition from pre-inflammatory to neutrophilic phase, testing whether preventing early alveolar macrophage permissiveness shifts that transition against the pathogen. Nutritional stresses such as iron limitation can be asked whether they hit harder in skin-restricted infections than in infections allowed to seed nodes, consistent with a placement-into-richer-niche pay-off.

Taken together, the biphasic lifestyle is cast as a staging strategy: use the macrophage environment briefly to solve early problems of transport, numbers, and incomplete arming, then abandon the intracellular lifestyle for extracellular, anti-phagocytic growth where the real population explosion occurs.

\section{Does phase two instigate macrophage bursting?}

A further question is whether maturation into the second, anti-phagocytic state actively causes the infected macrophage to die or rupture, releasing the now-armed bacteria into the extracellular space. Teaching summaries often put this bluntly---the bacteria replicate inside macrophages, then the macrophages burst and free bacilli spread---but the experimental literature is more careful. Temperature-induced virulence systems, especially the T3SS effector YopJ and its regulation by YopK, can drive macrophage death with solid experimental support. Evidence that this is how intracellular \textit{Y.\ pestis} exits after the first phase is weaker and incomplete. A coherent working model still links phase-two competence to programmed release without treating bursting as a solved fact.

\subsection{What is firmly known}

Macrophage death is a real part of \textit{Yersinia}--phagocyte biology. Pathogenic yersiniae induce death of macrophages and related cells. The best-characterised pathway is YopJ-dependent apoptosis. When a functional T3SS delivers YopJ into the host cytosol, YopJ acetylates components of MAPK and NF-$\kappa$B signalling, blocking pro-survival and pro-inflammatory transcription. Na\"ive macrophages are highly susceptible to this form of death. Classic work showed that YopJ is necessary for \textit{Yersinia}-triggered macrophage apoptosis, whereas several other Yops are not sufficient for that phenotype. YopJ-family activity is therefore a genetically supported link between the T3SS---a centrepiece of the phase-two anti-phagocytic package---and death of the very cell type used in phase one.

The outcome is not always quiet apoptosis. In activated macrophages, \textit{Yersinia} infection can be redirected toward caspase-1--dependent pyroptosis, a more inflammatory form of lysis. In lymph-node models of plague, necroptosis of infiltrated macrophages has been proposed to promote bacterial dissemination. Colloquial bursting may therefore correspond to different molecular programmes---apoptosis, pyroptosis, necroptosis, or secondary necrosis after apoptosis---depending on tissue, activation state, and timing.

Intracellular survival itself does not require the T3SS. Multiple studies indicate that \textit{Y.\ pestis} can survive and replicate inside macrophages in a manner that is largely independent of the pCD1/pYV-encoded T3SS. The YCV is remodelled by bacterial and host factors that are not identical to the Yop injection machinery used to paralyse extracellular phagocytosis. Phase one (living inside the macrophage) and phase two (killing or refusing macrophages via T3SS) are therefore experimentally separable. The T3SS is a leading candidate for how the bacterium later incapacitates or kills phagocytes when contact and injection become possible; it is not the main reason the bacterium survives intracellularly.

The exit step remains incompletely solved. Reviews of entry and resistance note that, after arming inside macrophages, \textit{Y.\ pestis} escapes to the extracellular environment, but that the mechanism of release is largely unknown, with apoptosis or necrosis in vitro offered as plausible correlates rather than as a closed mechanistic account. Textbooks assert bursting; primary literature still treats the precise exit route as an open problem.

\subsection{How phase two could cause release}

Even with that gap, a literature-grounded working model can answer the release question in the affirmative without elevating it to proven fact. Bacteria enter macrophages while relatively unarmed, especially after flea-temperature growth. Inside the YCV they replicate and, under host temperature and stress, induce the phase-two repertoire: F1, Psa, and the Ysc--Yop T3SS, including death-modulating effectors such as YopJ and regulators such as YopK. As bacterial load rises, or as bacteria gain access to the cytosolic face of host membranes after limited vacuolar compromise, host-cell damage, or re-encounter at the cell surface following partial release, T3SS injection becomes effective. YopJ, with related modulation of injection by YopK, then pushes the macrophage into apoptosis or another lethal programme; membrane integrity is lost; intracellular bacteria are liberated as an extracellular, anti-phagocytic population.

On this view, phase-two induction both prepares bacteria to resist the next phagocyte and, through part of the same package, destroys the current phagocyte, converting a staging niche into a release event. Macrophage death would then be an active, virulence-directed step rather than a passive side-effect of overcrowding.

\subsection{Supporting evidence}

YopJ apoptosis is genetically and molecularly linked to the T3SS. Loss of YopJ sharply reduces macrophage apoptosis after \textit{Yersinia} encounter. Because YopJ is delivered by the same secretion system that injects the anti-phagocytic Yops (YopH, YopE, YopT, YpkA), the death programme and the re-phagocytosis-resistance programme are co-packaged in the temperature-induced virulence plasmid system. Maturation toward phase two therefore includes molecular tools that can kill macrophages.

YopK ties macrophage apoptosis to disease-stage progression. In primary pneumonic plague, YopK modulates early macrophage apoptosis and has been shown to promote progression from the quieter early phase toward the later inflammatory stage of disease. That is direct experimental support for the idea that T3SS-controlled death of macrophages helps move infection through sequential stages, consistent with a biphasic bacterial strategy.

In vivo dissemination models also invoke phagocyte death. Work on bubonic spread has linked death programmes in infected or infiltrating mononuclear cells, including necroptosis in some settings, to release and onward movement of bacteria within and beyond lymph nodes. Teaching and review literature often summarises the sequence as intracellular replication followed by macrophage lysis and extracellular growth. Those summaries oversimplify mechanism, but they correctly reflect a consensus phenotype: successful plague involves a transition from cell-associated to free bacilli that coincides with destruction of host phagocytes.

Some data allow an intracellular contribution of YopJ. Although the textbook geometry of T3SS is injection by extracellular bacteria docked at the host surface, experimental systems with related yersiniae have reported that YopJ produced in the context of intracellular infection can contribute to higher levels of host-cell death. If similar routes operate for \textit{Y.\ pestis}---for example after YCV membrane damage or during incomplete phagocytosis---then bacteria maturing inside the cell could help trigger their own release without first becoming fully free.

\subsection{Limits and complications}

Efficient Yop injection is usually studied with extracellular bacteria contacting the plasma membrane. Bacteria sealed inside an intact YCV may not have the same access to inject YopJ into the cytosol. If so, phase-two maturation inside the vacuole would prepare anti-phagocytic competence for after release, while release itself might require a different trigger---host-initiated death, nutrient stress, physical overload, or a minority extracellular subpopulation that kills the cell from outside.

In some \textit{Y.\ pseudotuberculosis} macrophage models, a functional T3SS decreased survival of both bacterium and host after uptake, implying that premature or mis-timed expression of the system is not always beneficial inside the cell. That argues for tight temporal control: survive first with T3SS relatively quiet; deploy death and anti-phagocytosis later.

Not all phase-two factors are death factors. F1 capsule and the GTPase-targeting Yops that block uptake are not themselves the main apoptosis triggers. It is therefore imprecise to say that phase two as a whole causes bursting. More accurately, the temperature-induced virulence package that defines phase two includes both anti-phagocytic effectors and death-modulating effectors, and the latter are the candidates for active release.

Passive or host-driven exit also remains possible. Macrophages might die from cumulative stress, high intracellular bacterial burden, or immune signals independent of YopJ. In that case phase-two maturation would be concurrent with release rather than causal of it. Distinguishing correlation from causation requires time-resolved genetics---for example delayed induction of YopJ after established intracellular replication.

\subsection{Synthesis}

Relative to whether phase two instigates macrophage bursting, the literature supports a graded answer. Strongly supported is that phase-two-associated machinery (T3SS / YopJ, modulated by YopK) can and does kill macrophages; this is among the best-supported cell-biological effects of pathogenic \textit{Yersinia}. Moderately supported, at hypothesis level, is that such killing is a principal mechanism by which intracellular, matured \textit{Y.\ pestis} exits macrophages after phase one, liberating an extracellular population. That step is widely assumed and partially backed by disease-progression and dissemination studies, but it is not fully closed mechanistically. Weak or open remain simple physical bursting from overcrowding alone, and a single universal death pathway in all tissues and routes.

For the staging hypothesis developed above, the most useful formulation is therefore the following. After using the macrophage as a temporary niche for replication and virulence induction, \textit{Y.\ pestis} may deploy components of its phase-two programme---notably T3SS effectors that reprogramme host cell-death pathways---to destroy or rupture that niche, releasing bacteria that are now competent to resist re-phagocytosis. Anti-phagocytic Yops and F1 then maintain the extracellular state, while death-modulating effectors such as YopJ help create it. That assignment places bursting on a death-effector subset of the same induced system that defines phase two, not on phase two indiscriminately, and it leaves open the need for experiments that time T3SS competence to measured exit from the YCV in vivo.

\section*{Selected references}

\noindent
Ke, Y., Chen, Z.\ \& Yang, R.\ (2013). \textit{Yersinia pestis}: mechanisms of entry into and resistance to the host cell. \textit{Front.\ Cell.\ Infect.\ Microbiol.}\ 3, 106.\\
Du, Y., Rosqvist, R.\ \& Forsberg, \AA.\ (2002). Role of fraction 1 antigen of \textit{Yersinia pestis} in inhibition of phagocytosis. \textit{Infect.\ Immun.}\ 70, 1453--1460.\\
Pujol, C.\ \& Bliska, J.\ B.\ (2005). Turning \textit{Yersinia} pathogenesis outside in: subversion of macrophage function by intracellular yersiniae. \textit{Clin.\ Immunol.}\ 114, 216--226.\\
Pechous, R.\ D.\ et al.\ (2013). Early host cell targets of \textit{Yersinia pestis} during primary pneumonic plague. \textit{PLoS Pathog.}\ 9, e1003679.\\
Schwiesow, L.\ et al.\ (2016). \textit{Yersinia} type III secretion system master regulator LcrF. \textit{J.\ Bacteriol.}\ 198, 604--614.\\
Vadyvaloo, V.\ et al.\ (2010). Transit through the flea vector induces a pretransmission innate immunity resistance phenotype in \textit{Yersinia pestis}. \textit{PLoS Pathog.}\ 6, e1000783.\\
Sebbane, F.\ et al.\ (2008). The \textit{Yersinia pestis caf1M1A1} fimbrial capsule operon promotes transmission by flea bite in a mouse model of bubonic plague (including temperature dependence of F1: little or no capsule below about $35^\circ\mathrm{C}$ or in the flea). \textit{Infect.\ Immun.}\ 77, 1222--1229.\\
Lukaszewski, R.\ A.\ et al.\ (2005). Pathogenesis of \textit{Yersinia pestis} infection in BALB/c mice: effects on host macrophages and neutrophils. \textit{Infect.\ Immun.}\ 73, 7142--7150.\\
St.\ John, A.\ L.\ et al.\ (2014). S1P-dependent trafficking of intracellular \textit{Yersinia pestis} through lymph nodes establishes buboes and systemic infection. \textit{Immunity}\ 41, 440--450.\\
Bubeck, S.\ S.\ et al.\ (2007). Delayed inflammatory response to primary pneumonic plague occurs in both outbred and inbred mice. \textit{Infect.\ Immun.}\ 75, 697--705.\\
Price, S.\ L.\ et al.\ (2024). Microneedle array delivery of \textit{Yersinia pestis} recapitulates bubonic plague: rapid free drainage to lymph nodes can occur without obligatory intracellular carriage. \textit{iScience}.\\
Monack, D.\ M.\ et al.\ (1997). \textit{Yersinia} signals macrophages to undergo apoptosis and YopJ is necessary for this cell death. \textit{Proc.\ Natl.\ Acad.\ Sci.\ USA}\ 94, 10385--10390.\\
Bergsbaken, T.\ \& Cookson, B.\ T.\ (2007). Macrophage activation redirects \textit{Yersinia}-infected host cell death from apoptosis to caspase-1-dependent pyroptosis. \textit{PLoS Pathog.}\ 3, e161.\\
Peters, K.\ N.\ et al.\ (2013). Early apoptosis of macrophages modulated by injection of \textit{Yersinia pestis} YopK promotes progression of primary pneumonic plague. \textit{PLoS Pathog.}\ 9, e1003324.\\
Arifuzzaman, M.\ et al.\ (2018). Necroptosis of infiltrated macrophages drives \textit{Yersinia pestis} dissemination independent of IL-1$\beta$. Lymph-node models of phagocyte death and release.\\
Spinner, J.\ L.\ et al.\ (2014). \textit{Yersinia pestis} survival and replication within human neutrophils and implications for intracellular lifestyle relative to T3SS. \textit{J.\ Leukoc.\ Biol.}\ (and related work showing T3SS-independent intracellular survival phenotypes).

\end{document}
